\documentclass[aspectratio=169]{beamer}
% Shared preamble for ml/ slide decks.
% Each lecture .tex starts with:
%   \documentclass[aspectratio=169]{beamer}
%   % Shared preamble for ml/ slide decks.
% Each lecture .tex starts with:
%   \documentclass[aspectratio=169]{beamer}
%   % Shared preamble for ml/ slide decks.
% Each lecture .tex starts with:
%   \documentclass[aspectratio=169]{beamer}
%   \input{../preamble}
% Compile from inside the chapter folder so \input{../preamble} resolves.

\usetheme{default}
\usecolortheme{dove}

\setbeamertemplate{navigation symbols}{}
\setbeamertemplate{footline}{%
  \hfill{\large\insertframenumber\,/\,\inserttotalframenumber}\hspace{0.8em}\vspace{0.5em}%
}

% Project palette (existing, kept for backwards compatibility with deferred/*)
\definecolor{popblue}{RGB}{52, 101, 164}
\definecolor{sampred}{RGB}{204, 0, 0}
\definecolor{paramgreen}{RGB}{0, 140, 70}
\definecolor{warnred}{RGB}{180, 40, 40}
\definecolor{orange1}{RGB}{220, 120, 0}
\definecolor{violet1}{RGB}{120, 50, 160}
\definecolor{lightbg}{RGB}{245, 245, 250}

% Armenian flag palette (use these for any 3-color highlight; per CLAUDE.md).
% "Nice versions": slightly de-saturated from the official spec for slide use,
% so they don't burn out under projector colors. Originals kept as *_pure.
\definecolor{armred}{RGB}{200, 30, 40}      % nice red (official #D90012 muted)
\definecolor{armblue}{RGB}{30, 70, 160}     % nice blue (official #0033A0 muted)
\definecolor{armorange}{RGB}{230, 160, 30}  % nice orange (official #F2A800 muted)
\definecolor{armred_pure}{HTML}{D90012}     % official Armenian flag red
\definecolor{armblue_pure}{HTML}{0033A0}    % official Armenian flag blue
\definecolor{armorange_pure}{HTML}{F2A800}  % official Armenian flag orange

\setbeamercolor{frametitle}{fg=popblue}
\setbeamercolor{title}{fg=popblue}

\usepackage{pgfplots}
\usepackage{tikz}
\usetikzlibrary{shapes, arrows.meta, positioning, calc, decorations.pathreplacing, patterns, matrix}
\pgfplotsset{compat=1.18}
\usepackage{amsmath, amssymb, amsfonts}
\usepackage{bm}
\usepackage{dsfont}
\usepackage{array}
\usepackage{pifont}
\usepackage{colortbl}
\usepackage{fontenc}

% Callout boxes. Purely additive - existing decks are unaffected.
% Usage: \begin{tcolorbox}[colback=armblue!8, colframe=armblue, boxrule=0.6pt, arc=2pt,
%                          title=\textbf{Key idea}] ... \end{tcolorbox}
% Palette convention (ml/SLIDE_STYLE.md): armblue=key idea, armred=trap/warning,
% paramgreen=takeaway or "Next", armorange=watch-out.
\usepackage[most]{tcolorbox}

% Cleaner table rules (\toprule, \midrule, \bottomrule). Also purely additive.
\usepackage{booktabs}

% \cancel{...} for struck-through terms in derivations (added 2026-08-11 for the
% importance-sampling frame in ch11_rl, where the dynamics cancel). Purely additive.
\usepackage{cancel}

% Minimal ML notation macros (subset of slds-lmu latex-math/basic-ml.tex).
% Kept here so we don't have to inline-expand every \thetav, \xv, etc.
% when porting upstream slides.
\newcommand{\R}{\mathds{R}}
\newcommand{\N}{\mathds{N}}
\newcommand{\E}{\mathds{E}}
\newcommand{\Xspace}{\mathcal{X}}
\newcommand{\Yspace}{\mathcal{Y}}
\newcommand{\Hspace}{\mathcal{H}}
\newcommand{\D}{\mathcal{D}}
\newcommand{\Dtrain}{\mathcal{D}_{\text{train}}}
\newcommand{\Dtest}{\mathcal{D}_{\text{test}}}
\newcommand{\xv}{\mathbf{x}}
\newcommand{\yv}{\mathbf{y}}
\newcommand{\Xmat}{\mathbf{X}}
\newcommand{\id}{\mathbf{I}}
\newcommand{\thetav}{\bm{\theta}}
\newcommand{\thetah}{\hat{\bm{\theta}}}
\newcommand{\lamv}{\bm{\lambda}}
\newcommand{\fx}{f(\xv)}
\newcommand{\fxt}{f(\xv \,|\, \bm{\theta})}
\newcommand{\fxi}{f\!\left(\xv^{(i)}\right)}
\newcommand{\fh}{\hat{f}}
\newcommand{\fxh}{\hat{f}(\xv)}
\newcommand{\yh}{\hat{y}}
\newcommand{\yih}{\hat{y}^{(i)}}
% \xi clashes with the Greek letter \xi; write \xv^{(i)} inline instead.
\newcommand{\yi}[1][i]{y^{(#1)}}
\newcommand{\xyi}[1][i]{\left(\xv^{(#1)},\, y^{(#1)}\right)}
\newcommand{\thx}{\bm{\theta}^\top \xv}
\newcommand{\sumin}{\sum_{i=1}^{n}}
\newcommand{\sumjp}{\sum_{j=1}^{p}}
\newcommand{\Lxy}{L\!\left(y, \fx\right)}
\newcommand{\Lxyi}{L\!\left(\yi, \fxi\right)}
\newcommand{\risk}{\mathcal{R}}
\newcommand{\riske}{\mathcal{R}_{\text{emp}}}
\newcommand{\risket}{\mathcal{R}_{\text{emp}}(\bm{\theta})}
\newcommand{\riskt}{\mathcal{R}(\bm{\theta})}
\newcommand{\riskr}{\mathcal{R}_{\text{reg}}}
\newcommand{\argmin}{\mathop{\mathrm{arg\,min}}}
\newcommand{\argmax}{\mathop{\mathrm{arg\,max}}}
\newcommand{\eps}{\varepsilon}

% Compile from inside the chapter folder so % Shared preamble for ml/ slide decks.
% Each lecture .tex starts with:
%   \documentclass[aspectratio=169]{beamer}
%   \input{../preamble}
% Compile from inside the chapter folder so \input{../preamble} resolves.

\usetheme{default}
\usecolortheme{dove}

\setbeamertemplate{navigation symbols}{}
\setbeamertemplate{footline}{%
  \hfill{\large\insertframenumber\,/\,\inserttotalframenumber}\hspace{0.8em}\vspace{0.5em}%
}

% Project palette (existing, kept for backwards compatibility with deferred/*)
\definecolor{popblue}{RGB}{52, 101, 164}
\definecolor{sampred}{RGB}{204, 0, 0}
\definecolor{paramgreen}{RGB}{0, 140, 70}
\definecolor{warnred}{RGB}{180, 40, 40}
\definecolor{orange1}{RGB}{220, 120, 0}
\definecolor{violet1}{RGB}{120, 50, 160}
\definecolor{lightbg}{RGB}{245, 245, 250}

% Armenian flag palette (use these for any 3-color highlight; per CLAUDE.md).
% "Nice versions": slightly de-saturated from the official spec for slide use,
% so they don't burn out under projector colors. Originals kept as *_pure.
\definecolor{armred}{RGB}{200, 30, 40}      % nice red (official #D90012 muted)
\definecolor{armblue}{RGB}{30, 70, 160}     % nice blue (official #0033A0 muted)
\definecolor{armorange}{RGB}{230, 160, 30}  % nice orange (official #F2A800 muted)
\definecolor{armred_pure}{HTML}{D90012}     % official Armenian flag red
\definecolor{armblue_pure}{HTML}{0033A0}    % official Armenian flag blue
\definecolor{armorange_pure}{HTML}{F2A800}  % official Armenian flag orange

\setbeamercolor{frametitle}{fg=popblue}
\setbeamercolor{title}{fg=popblue}

\usepackage{pgfplots}
\usepackage{tikz}
\usetikzlibrary{shapes, arrows.meta, positioning, calc, decorations.pathreplacing, patterns, matrix}
\pgfplotsset{compat=1.18}
\usepackage{amsmath, amssymb, amsfonts}
\usepackage{bm}
\usepackage{dsfont}
\usepackage{array}
\usepackage{pifont}
\usepackage{colortbl}
\usepackage{fontenc}

% Callout boxes. Purely additive - existing decks are unaffected.
% Usage: \begin{tcolorbox}[colback=armblue!8, colframe=armblue, boxrule=0.6pt, arc=2pt,
%                          title=\textbf{Key idea}] ... \end{tcolorbox}
% Palette convention (ml/SLIDE_STYLE.md): armblue=key idea, armred=trap/warning,
% paramgreen=takeaway or "Next", armorange=watch-out.
\usepackage[most]{tcolorbox}

% Cleaner table rules (\toprule, \midrule, \bottomrule). Also purely additive.
\usepackage{booktabs}

% \cancel{...} for struck-through terms in derivations (added 2026-08-11 for the
% importance-sampling frame in ch11_rl, where the dynamics cancel). Purely additive.
\usepackage{cancel}

% Minimal ML notation macros (subset of slds-lmu latex-math/basic-ml.tex).
% Kept here so we don't have to inline-expand every \thetav, \xv, etc.
% when porting upstream slides.
\newcommand{\R}{\mathds{R}}
\newcommand{\N}{\mathds{N}}
\newcommand{\E}{\mathds{E}}
\newcommand{\Xspace}{\mathcal{X}}
\newcommand{\Yspace}{\mathcal{Y}}
\newcommand{\Hspace}{\mathcal{H}}
\newcommand{\D}{\mathcal{D}}
\newcommand{\Dtrain}{\mathcal{D}_{\text{train}}}
\newcommand{\Dtest}{\mathcal{D}_{\text{test}}}
\newcommand{\xv}{\mathbf{x}}
\newcommand{\yv}{\mathbf{y}}
\newcommand{\Xmat}{\mathbf{X}}
\newcommand{\id}{\mathbf{I}}
\newcommand{\thetav}{\bm{\theta}}
\newcommand{\thetah}{\hat{\bm{\theta}}}
\newcommand{\lamv}{\bm{\lambda}}
\newcommand{\fx}{f(\xv)}
\newcommand{\fxt}{f(\xv \,|\, \bm{\theta})}
\newcommand{\fxi}{f\!\left(\xv^{(i)}\right)}
\newcommand{\fh}{\hat{f}}
\newcommand{\fxh}{\hat{f}(\xv)}
\newcommand{\yh}{\hat{y}}
\newcommand{\yih}{\hat{y}^{(i)}}
% \xi clashes with the Greek letter \xi; write \xv^{(i)} inline instead.
\newcommand{\yi}[1][i]{y^{(#1)}}
\newcommand{\xyi}[1][i]{\left(\xv^{(#1)},\, y^{(#1)}\right)}
\newcommand{\thx}{\bm{\theta}^\top \xv}
\newcommand{\sumin}{\sum_{i=1}^{n}}
\newcommand{\sumjp}{\sum_{j=1}^{p}}
\newcommand{\Lxy}{L\!\left(y, \fx\right)}
\newcommand{\Lxyi}{L\!\left(\yi, \fxi\right)}
\newcommand{\risk}{\mathcal{R}}
\newcommand{\riske}{\mathcal{R}_{\text{emp}}}
\newcommand{\risket}{\mathcal{R}_{\text{emp}}(\bm{\theta})}
\newcommand{\riskt}{\mathcal{R}(\bm{\theta})}
\newcommand{\riskr}{\mathcal{R}_{\text{reg}}}
\newcommand{\argmin}{\mathop{\mathrm{arg\,min}}}
\newcommand{\argmax}{\mathop{\mathrm{arg\,max}}}
\newcommand{\eps}{\varepsilon}
 resolves.

\usetheme{default}
\usecolortheme{dove}

\setbeamertemplate{navigation symbols}{}
\setbeamertemplate{footline}{%
  \hfill{\large\insertframenumber\,/\,\inserttotalframenumber}\hspace{0.8em}\vspace{0.5em}%
}

% Project palette (existing, kept for backwards compatibility with deferred/*)
\definecolor{popblue}{RGB}{52, 101, 164}
\definecolor{sampred}{RGB}{204, 0, 0}
\definecolor{paramgreen}{RGB}{0, 140, 70}
\definecolor{warnred}{RGB}{180, 40, 40}
\definecolor{orange1}{RGB}{220, 120, 0}
\definecolor{violet1}{RGB}{120, 50, 160}
\definecolor{lightbg}{RGB}{245, 245, 250}

% Armenian flag palette (use these for any 3-color highlight; per CLAUDE.md).
% "Nice versions": slightly de-saturated from the official spec for slide use,
% so they don't burn out under projector colors. Originals kept as *_pure.
\definecolor{armred}{RGB}{200, 30, 40}      % nice red (official #D90012 muted)
\definecolor{armblue}{RGB}{30, 70, 160}     % nice blue (official #0033A0 muted)
\definecolor{armorange}{RGB}{230, 160, 30}  % nice orange (official #F2A800 muted)
\definecolor{armred_pure}{HTML}{D90012}     % official Armenian flag red
\definecolor{armblue_pure}{HTML}{0033A0}    % official Armenian flag blue
\definecolor{armorange_pure}{HTML}{F2A800}  % official Armenian flag orange

\setbeamercolor{frametitle}{fg=popblue}
\setbeamercolor{title}{fg=popblue}

\usepackage{pgfplots}
\usepackage{tikz}
\usetikzlibrary{shapes, arrows.meta, positioning, calc, decorations.pathreplacing, patterns, matrix}
\pgfplotsset{compat=1.18}
\usepackage{amsmath, amssymb, amsfonts}
\usepackage{bm}
\usepackage{dsfont}
\usepackage{array}
\usepackage{pifont}
\usepackage{colortbl}
\usepackage{fontenc}

% Callout boxes. Purely additive - existing decks are unaffected.
% Usage: \begin{tcolorbox}[colback=armblue!8, colframe=armblue, boxrule=0.6pt, arc=2pt,
%                          title=\textbf{Key idea}] ... \end{tcolorbox}
% Palette convention (ml/SLIDE_STYLE.md): armblue=key idea, armred=trap/warning,
% paramgreen=takeaway or "Next", armorange=watch-out.
\usepackage[most]{tcolorbox}

% Cleaner table rules (\toprule, \midrule, \bottomrule). Also purely additive.
\usepackage{booktabs}

% \cancel{...} for struck-through terms in derivations (added 2026-08-11 for the
% importance-sampling frame in ch11_rl, where the dynamics cancel). Purely additive.
\usepackage{cancel}

% Minimal ML notation macros (subset of slds-lmu latex-math/basic-ml.tex).
% Kept here so we don't have to inline-expand every \thetav, \xv, etc.
% when porting upstream slides.
\newcommand{\R}{\mathds{R}}
\newcommand{\N}{\mathds{N}}
\newcommand{\E}{\mathds{E}}
\newcommand{\Xspace}{\mathcal{X}}
\newcommand{\Yspace}{\mathcal{Y}}
\newcommand{\Hspace}{\mathcal{H}}
\newcommand{\D}{\mathcal{D}}
\newcommand{\Dtrain}{\mathcal{D}_{\text{train}}}
\newcommand{\Dtest}{\mathcal{D}_{\text{test}}}
\newcommand{\xv}{\mathbf{x}}
\newcommand{\yv}{\mathbf{y}}
\newcommand{\Xmat}{\mathbf{X}}
\newcommand{\id}{\mathbf{I}}
\newcommand{\thetav}{\bm{\theta}}
\newcommand{\thetah}{\hat{\bm{\theta}}}
\newcommand{\lamv}{\bm{\lambda}}
\newcommand{\fx}{f(\xv)}
\newcommand{\fxt}{f(\xv \,|\, \bm{\theta})}
\newcommand{\fxi}{f\!\left(\xv^{(i)}\right)}
\newcommand{\fh}{\hat{f}}
\newcommand{\fxh}{\hat{f}(\xv)}
\newcommand{\yh}{\hat{y}}
\newcommand{\yih}{\hat{y}^{(i)}}
% \xi clashes with the Greek letter \xi; write \xv^{(i)} inline instead.
\newcommand{\yi}[1][i]{y^{(#1)}}
\newcommand{\xyi}[1][i]{\left(\xv^{(#1)},\, y^{(#1)}\right)}
\newcommand{\thx}{\bm{\theta}^\top \xv}
\newcommand{\sumin}{\sum_{i=1}^{n}}
\newcommand{\sumjp}{\sum_{j=1}^{p}}
\newcommand{\Lxy}{L\!\left(y, \fx\right)}
\newcommand{\Lxyi}{L\!\left(\yi, \fxi\right)}
\newcommand{\risk}{\mathcal{R}}
\newcommand{\riske}{\mathcal{R}_{\text{emp}}}
\newcommand{\risket}{\mathcal{R}_{\text{emp}}(\bm{\theta})}
\newcommand{\riskt}{\mathcal{R}(\bm{\theta})}
\newcommand{\riskr}{\mathcal{R}_{\text{reg}}}
\newcommand{\argmin}{\mathop{\mathrm{arg\,min}}}
\newcommand{\argmax}{\mathop{\mathrm{arg\,max}}}
\newcommand{\eps}{\varepsilon}

% Compile from inside the chapter folder so % Shared preamble for ml/ slide decks.
% Each lecture .tex starts with:
%   \documentclass[aspectratio=169]{beamer}
%   % Shared preamble for ml/ slide decks.
% Each lecture .tex starts with:
%   \documentclass[aspectratio=169]{beamer}
%   \input{../preamble}
% Compile from inside the chapter folder so \input{../preamble} resolves.

\usetheme{default}
\usecolortheme{dove}

\setbeamertemplate{navigation symbols}{}
\setbeamertemplate{footline}{%
  \hfill{\large\insertframenumber\,/\,\inserttotalframenumber}\hspace{0.8em}\vspace{0.5em}%
}

% Project palette (existing, kept for backwards compatibility with deferred/*)
\definecolor{popblue}{RGB}{52, 101, 164}
\definecolor{sampred}{RGB}{204, 0, 0}
\definecolor{paramgreen}{RGB}{0, 140, 70}
\definecolor{warnred}{RGB}{180, 40, 40}
\definecolor{orange1}{RGB}{220, 120, 0}
\definecolor{violet1}{RGB}{120, 50, 160}
\definecolor{lightbg}{RGB}{245, 245, 250}

% Armenian flag palette (use these for any 3-color highlight; per CLAUDE.md).
% "Nice versions": slightly de-saturated from the official spec for slide use,
% so they don't burn out under projector colors. Originals kept as *_pure.
\definecolor{armred}{RGB}{200, 30, 40}      % nice red (official #D90012 muted)
\definecolor{armblue}{RGB}{30, 70, 160}     % nice blue (official #0033A0 muted)
\definecolor{armorange}{RGB}{230, 160, 30}  % nice orange (official #F2A800 muted)
\definecolor{armred_pure}{HTML}{D90012}     % official Armenian flag red
\definecolor{armblue_pure}{HTML}{0033A0}    % official Armenian flag blue
\definecolor{armorange_pure}{HTML}{F2A800}  % official Armenian flag orange

\setbeamercolor{frametitle}{fg=popblue}
\setbeamercolor{title}{fg=popblue}

\usepackage{pgfplots}
\usepackage{tikz}
\usetikzlibrary{shapes, arrows.meta, positioning, calc, decorations.pathreplacing, patterns, matrix}
\pgfplotsset{compat=1.18}
\usepackage{amsmath, amssymb, amsfonts}
\usepackage{bm}
\usepackage{dsfont}
\usepackage{array}
\usepackage{pifont}
\usepackage{colortbl}
\usepackage{fontenc}

% Callout boxes. Purely additive - existing decks are unaffected.
% Usage: \begin{tcolorbox}[colback=armblue!8, colframe=armblue, boxrule=0.6pt, arc=2pt,
%                          title=\textbf{Key idea}] ... \end{tcolorbox}
% Palette convention (ml/SLIDE_STYLE.md): armblue=key idea, armred=trap/warning,
% paramgreen=takeaway or "Next", armorange=watch-out.
\usepackage[most]{tcolorbox}

% Cleaner table rules (\toprule, \midrule, \bottomrule). Also purely additive.
\usepackage{booktabs}

% \cancel{...} for struck-through terms in derivations (added 2026-08-11 for the
% importance-sampling frame in ch11_rl, where the dynamics cancel). Purely additive.
\usepackage{cancel}

% Minimal ML notation macros (subset of slds-lmu latex-math/basic-ml.tex).
% Kept here so we don't have to inline-expand every \thetav, \xv, etc.
% when porting upstream slides.
\newcommand{\R}{\mathds{R}}
\newcommand{\N}{\mathds{N}}
\newcommand{\E}{\mathds{E}}
\newcommand{\Xspace}{\mathcal{X}}
\newcommand{\Yspace}{\mathcal{Y}}
\newcommand{\Hspace}{\mathcal{H}}
\newcommand{\D}{\mathcal{D}}
\newcommand{\Dtrain}{\mathcal{D}_{\text{train}}}
\newcommand{\Dtest}{\mathcal{D}_{\text{test}}}
\newcommand{\xv}{\mathbf{x}}
\newcommand{\yv}{\mathbf{y}}
\newcommand{\Xmat}{\mathbf{X}}
\newcommand{\id}{\mathbf{I}}
\newcommand{\thetav}{\bm{\theta}}
\newcommand{\thetah}{\hat{\bm{\theta}}}
\newcommand{\lamv}{\bm{\lambda}}
\newcommand{\fx}{f(\xv)}
\newcommand{\fxt}{f(\xv \,|\, \bm{\theta})}
\newcommand{\fxi}{f\!\left(\xv^{(i)}\right)}
\newcommand{\fh}{\hat{f}}
\newcommand{\fxh}{\hat{f}(\xv)}
\newcommand{\yh}{\hat{y}}
\newcommand{\yih}{\hat{y}^{(i)}}
% \xi clashes with the Greek letter \xi; write \xv^{(i)} inline instead.
\newcommand{\yi}[1][i]{y^{(#1)}}
\newcommand{\xyi}[1][i]{\left(\xv^{(#1)},\, y^{(#1)}\right)}
\newcommand{\thx}{\bm{\theta}^\top \xv}
\newcommand{\sumin}{\sum_{i=1}^{n}}
\newcommand{\sumjp}{\sum_{j=1}^{p}}
\newcommand{\Lxy}{L\!\left(y, \fx\right)}
\newcommand{\Lxyi}{L\!\left(\yi, \fxi\right)}
\newcommand{\risk}{\mathcal{R}}
\newcommand{\riske}{\mathcal{R}_{\text{emp}}}
\newcommand{\risket}{\mathcal{R}_{\text{emp}}(\bm{\theta})}
\newcommand{\riskt}{\mathcal{R}(\bm{\theta})}
\newcommand{\riskr}{\mathcal{R}_{\text{reg}}}
\newcommand{\argmin}{\mathop{\mathrm{arg\,min}}}
\newcommand{\argmax}{\mathop{\mathrm{arg\,max}}}
\newcommand{\eps}{\varepsilon}

% Compile from inside the chapter folder so % Shared preamble for ml/ slide decks.
% Each lecture .tex starts with:
%   \documentclass[aspectratio=169]{beamer}
%   \input{../preamble}
% Compile from inside the chapter folder so \input{../preamble} resolves.

\usetheme{default}
\usecolortheme{dove}

\setbeamertemplate{navigation symbols}{}
\setbeamertemplate{footline}{%
  \hfill{\large\insertframenumber\,/\,\inserttotalframenumber}\hspace{0.8em}\vspace{0.5em}%
}

% Project palette (existing, kept for backwards compatibility with deferred/*)
\definecolor{popblue}{RGB}{52, 101, 164}
\definecolor{sampred}{RGB}{204, 0, 0}
\definecolor{paramgreen}{RGB}{0, 140, 70}
\definecolor{warnred}{RGB}{180, 40, 40}
\definecolor{orange1}{RGB}{220, 120, 0}
\definecolor{violet1}{RGB}{120, 50, 160}
\definecolor{lightbg}{RGB}{245, 245, 250}

% Armenian flag palette (use these for any 3-color highlight; per CLAUDE.md).
% "Nice versions": slightly de-saturated from the official spec for slide use,
% so they don't burn out under projector colors. Originals kept as *_pure.
\definecolor{armred}{RGB}{200, 30, 40}      % nice red (official #D90012 muted)
\definecolor{armblue}{RGB}{30, 70, 160}     % nice blue (official #0033A0 muted)
\definecolor{armorange}{RGB}{230, 160, 30}  % nice orange (official #F2A800 muted)
\definecolor{armred_pure}{HTML}{D90012}     % official Armenian flag red
\definecolor{armblue_pure}{HTML}{0033A0}    % official Armenian flag blue
\definecolor{armorange_pure}{HTML}{F2A800}  % official Armenian flag orange

\setbeamercolor{frametitle}{fg=popblue}
\setbeamercolor{title}{fg=popblue}

\usepackage{pgfplots}
\usepackage{tikz}
\usetikzlibrary{shapes, arrows.meta, positioning, calc, decorations.pathreplacing, patterns, matrix}
\pgfplotsset{compat=1.18}
\usepackage{amsmath, amssymb, amsfonts}
\usepackage{bm}
\usepackage{dsfont}
\usepackage{array}
\usepackage{pifont}
\usepackage{colortbl}
\usepackage{fontenc}

% Callout boxes. Purely additive - existing decks are unaffected.
% Usage: \begin{tcolorbox}[colback=armblue!8, colframe=armblue, boxrule=0.6pt, arc=2pt,
%                          title=\textbf{Key idea}] ... \end{tcolorbox}
% Palette convention (ml/SLIDE_STYLE.md): armblue=key idea, armred=trap/warning,
% paramgreen=takeaway or "Next", armorange=watch-out.
\usepackage[most]{tcolorbox}

% Cleaner table rules (\toprule, \midrule, \bottomrule). Also purely additive.
\usepackage{booktabs}

% \cancel{...} for struck-through terms in derivations (added 2026-08-11 for the
% importance-sampling frame in ch11_rl, where the dynamics cancel). Purely additive.
\usepackage{cancel}

% Minimal ML notation macros (subset of slds-lmu latex-math/basic-ml.tex).
% Kept here so we don't have to inline-expand every \thetav, \xv, etc.
% when porting upstream slides.
\newcommand{\R}{\mathds{R}}
\newcommand{\N}{\mathds{N}}
\newcommand{\E}{\mathds{E}}
\newcommand{\Xspace}{\mathcal{X}}
\newcommand{\Yspace}{\mathcal{Y}}
\newcommand{\Hspace}{\mathcal{H}}
\newcommand{\D}{\mathcal{D}}
\newcommand{\Dtrain}{\mathcal{D}_{\text{train}}}
\newcommand{\Dtest}{\mathcal{D}_{\text{test}}}
\newcommand{\xv}{\mathbf{x}}
\newcommand{\yv}{\mathbf{y}}
\newcommand{\Xmat}{\mathbf{X}}
\newcommand{\id}{\mathbf{I}}
\newcommand{\thetav}{\bm{\theta}}
\newcommand{\thetah}{\hat{\bm{\theta}}}
\newcommand{\lamv}{\bm{\lambda}}
\newcommand{\fx}{f(\xv)}
\newcommand{\fxt}{f(\xv \,|\, \bm{\theta})}
\newcommand{\fxi}{f\!\left(\xv^{(i)}\right)}
\newcommand{\fh}{\hat{f}}
\newcommand{\fxh}{\hat{f}(\xv)}
\newcommand{\yh}{\hat{y}}
\newcommand{\yih}{\hat{y}^{(i)}}
% \xi clashes with the Greek letter \xi; write \xv^{(i)} inline instead.
\newcommand{\yi}[1][i]{y^{(#1)}}
\newcommand{\xyi}[1][i]{\left(\xv^{(#1)},\, y^{(#1)}\right)}
\newcommand{\thx}{\bm{\theta}^\top \xv}
\newcommand{\sumin}{\sum_{i=1}^{n}}
\newcommand{\sumjp}{\sum_{j=1}^{p}}
\newcommand{\Lxy}{L\!\left(y, \fx\right)}
\newcommand{\Lxyi}{L\!\left(\yi, \fxi\right)}
\newcommand{\risk}{\mathcal{R}}
\newcommand{\riske}{\mathcal{R}_{\text{emp}}}
\newcommand{\risket}{\mathcal{R}_{\text{emp}}(\bm{\theta})}
\newcommand{\riskt}{\mathcal{R}(\bm{\theta})}
\newcommand{\riskr}{\mathcal{R}_{\text{reg}}}
\newcommand{\argmin}{\mathop{\mathrm{arg\,min}}}
\newcommand{\argmax}{\mathop{\mathrm{arg\,max}}}
\newcommand{\eps}{\varepsilon}
 resolves.

\usetheme{default}
\usecolortheme{dove}

\setbeamertemplate{navigation symbols}{}
\setbeamertemplate{footline}{%
  \hfill{\large\insertframenumber\,/\,\inserttotalframenumber}\hspace{0.8em}\vspace{0.5em}%
}

% Project palette (existing, kept for backwards compatibility with deferred/*)
\definecolor{popblue}{RGB}{52, 101, 164}
\definecolor{sampred}{RGB}{204, 0, 0}
\definecolor{paramgreen}{RGB}{0, 140, 70}
\definecolor{warnred}{RGB}{180, 40, 40}
\definecolor{orange1}{RGB}{220, 120, 0}
\definecolor{violet1}{RGB}{120, 50, 160}
\definecolor{lightbg}{RGB}{245, 245, 250}

% Armenian flag palette (use these for any 3-color highlight; per CLAUDE.md).
% "Nice versions": slightly de-saturated from the official spec for slide use,
% so they don't burn out under projector colors. Originals kept as *_pure.
\definecolor{armred}{RGB}{200, 30, 40}      % nice red (official #D90012 muted)
\definecolor{armblue}{RGB}{30, 70, 160}     % nice blue (official #0033A0 muted)
\definecolor{armorange}{RGB}{230, 160, 30}  % nice orange (official #F2A800 muted)
\definecolor{armred_pure}{HTML}{D90012}     % official Armenian flag red
\definecolor{armblue_pure}{HTML}{0033A0}    % official Armenian flag blue
\definecolor{armorange_pure}{HTML}{F2A800}  % official Armenian flag orange

\setbeamercolor{frametitle}{fg=popblue}
\setbeamercolor{title}{fg=popblue}

\usepackage{pgfplots}
\usepackage{tikz}
\usetikzlibrary{shapes, arrows.meta, positioning, calc, decorations.pathreplacing, patterns, matrix}
\pgfplotsset{compat=1.18}
\usepackage{amsmath, amssymb, amsfonts}
\usepackage{bm}
\usepackage{dsfont}
\usepackage{array}
\usepackage{pifont}
\usepackage{colortbl}
\usepackage{fontenc}

% Callout boxes. Purely additive - existing decks are unaffected.
% Usage: \begin{tcolorbox}[colback=armblue!8, colframe=armblue, boxrule=0.6pt, arc=2pt,
%                          title=\textbf{Key idea}] ... \end{tcolorbox}
% Palette convention (ml/SLIDE_STYLE.md): armblue=key idea, armred=trap/warning,
% paramgreen=takeaway or "Next", armorange=watch-out.
\usepackage[most]{tcolorbox}

% Cleaner table rules (\toprule, \midrule, \bottomrule). Also purely additive.
\usepackage{booktabs}

% \cancel{...} for struck-through terms in derivations (added 2026-08-11 for the
% importance-sampling frame in ch11_rl, where the dynamics cancel). Purely additive.
\usepackage{cancel}

% Minimal ML notation macros (subset of slds-lmu latex-math/basic-ml.tex).
% Kept here so we don't have to inline-expand every \thetav, \xv, etc.
% when porting upstream slides.
\newcommand{\R}{\mathds{R}}
\newcommand{\N}{\mathds{N}}
\newcommand{\E}{\mathds{E}}
\newcommand{\Xspace}{\mathcal{X}}
\newcommand{\Yspace}{\mathcal{Y}}
\newcommand{\Hspace}{\mathcal{H}}
\newcommand{\D}{\mathcal{D}}
\newcommand{\Dtrain}{\mathcal{D}_{\text{train}}}
\newcommand{\Dtest}{\mathcal{D}_{\text{test}}}
\newcommand{\xv}{\mathbf{x}}
\newcommand{\yv}{\mathbf{y}}
\newcommand{\Xmat}{\mathbf{X}}
\newcommand{\id}{\mathbf{I}}
\newcommand{\thetav}{\bm{\theta}}
\newcommand{\thetah}{\hat{\bm{\theta}}}
\newcommand{\lamv}{\bm{\lambda}}
\newcommand{\fx}{f(\xv)}
\newcommand{\fxt}{f(\xv \,|\, \bm{\theta})}
\newcommand{\fxi}{f\!\left(\xv^{(i)}\right)}
\newcommand{\fh}{\hat{f}}
\newcommand{\fxh}{\hat{f}(\xv)}
\newcommand{\yh}{\hat{y}}
\newcommand{\yih}{\hat{y}^{(i)}}
% \xi clashes with the Greek letter \xi; write \xv^{(i)} inline instead.
\newcommand{\yi}[1][i]{y^{(#1)}}
\newcommand{\xyi}[1][i]{\left(\xv^{(#1)},\, y^{(#1)}\right)}
\newcommand{\thx}{\bm{\theta}^\top \xv}
\newcommand{\sumin}{\sum_{i=1}^{n}}
\newcommand{\sumjp}{\sum_{j=1}^{p}}
\newcommand{\Lxy}{L\!\left(y, \fx\right)}
\newcommand{\Lxyi}{L\!\left(\yi, \fxi\right)}
\newcommand{\risk}{\mathcal{R}}
\newcommand{\riske}{\mathcal{R}_{\text{emp}}}
\newcommand{\risket}{\mathcal{R}_{\text{emp}}(\bm{\theta})}
\newcommand{\riskt}{\mathcal{R}(\bm{\theta})}
\newcommand{\riskr}{\mathcal{R}_{\text{reg}}}
\newcommand{\argmin}{\mathop{\mathrm{arg\,min}}}
\newcommand{\argmax}{\mathop{\mathrm{arg\,max}}}
\newcommand{\eps}{\varepsilon}
 resolves.

\usetheme{default}
\usecolortheme{dove}

\setbeamertemplate{navigation symbols}{}
\setbeamertemplate{footline}{%
  \hfill{\large\insertframenumber\,/\,\inserttotalframenumber}\hspace{0.8em}\vspace{0.5em}%
}

% Project palette (existing, kept for backwards compatibility with deferred/*)
\definecolor{popblue}{RGB}{52, 101, 164}
\definecolor{sampred}{RGB}{204, 0, 0}
\definecolor{paramgreen}{RGB}{0, 140, 70}
\definecolor{warnred}{RGB}{180, 40, 40}
\definecolor{orange1}{RGB}{220, 120, 0}
\definecolor{violet1}{RGB}{120, 50, 160}
\definecolor{lightbg}{RGB}{245, 245, 250}

% Armenian flag palette (use these for any 3-color highlight; per CLAUDE.md).
% "Nice versions": slightly de-saturated from the official spec for slide use,
% so they don't burn out under projector colors. Originals kept as *_pure.
\definecolor{armred}{RGB}{200, 30, 40}      % nice red (official #D90012 muted)
\definecolor{armblue}{RGB}{30, 70, 160}     % nice blue (official #0033A0 muted)
\definecolor{armorange}{RGB}{230, 160, 30}  % nice orange (official #F2A800 muted)
\definecolor{armred_pure}{HTML}{D90012}     % official Armenian flag red
\definecolor{armblue_pure}{HTML}{0033A0}    % official Armenian flag blue
\definecolor{armorange_pure}{HTML}{F2A800}  % official Armenian flag orange

\setbeamercolor{frametitle}{fg=popblue}
\setbeamercolor{title}{fg=popblue}

\usepackage{pgfplots}
\usepackage{tikz}
\usetikzlibrary{shapes, arrows.meta, positioning, calc, decorations.pathreplacing, patterns, matrix}
\pgfplotsset{compat=1.18}
\usepackage{amsmath, amssymb, amsfonts}
\usepackage{bm}
\usepackage{dsfont}
\usepackage{array}
\usepackage{pifont}
\usepackage{colortbl}
\usepackage{fontenc}

% Callout boxes. Purely additive - existing decks are unaffected.
% Usage: \begin{tcolorbox}[colback=armblue!8, colframe=armblue, boxrule=0.6pt, arc=2pt,
%                          title=\textbf{Key idea}] ... \end{tcolorbox}
% Palette convention (ml/SLIDE_STYLE.md): armblue=key idea, armred=trap/warning,
% paramgreen=takeaway or "Next", armorange=watch-out.
\usepackage[most]{tcolorbox}

% Cleaner table rules (\toprule, \midrule, \bottomrule). Also purely additive.
\usepackage{booktabs}

% \cancel{...} for struck-through terms in derivations (added 2026-08-11 for the
% importance-sampling frame in ch11_rl, where the dynamics cancel). Purely additive.
\usepackage{cancel}

% Minimal ML notation macros (subset of slds-lmu latex-math/basic-ml.tex).
% Kept here so we don't have to inline-expand every \thetav, \xv, etc.
% when porting upstream slides.
\newcommand{\R}{\mathds{R}}
\newcommand{\N}{\mathds{N}}
\newcommand{\E}{\mathds{E}}
\newcommand{\Xspace}{\mathcal{X}}
\newcommand{\Yspace}{\mathcal{Y}}
\newcommand{\Hspace}{\mathcal{H}}
\newcommand{\D}{\mathcal{D}}
\newcommand{\Dtrain}{\mathcal{D}_{\text{train}}}
\newcommand{\Dtest}{\mathcal{D}_{\text{test}}}
\newcommand{\xv}{\mathbf{x}}
\newcommand{\yv}{\mathbf{y}}
\newcommand{\Xmat}{\mathbf{X}}
\newcommand{\id}{\mathbf{I}}
\newcommand{\thetav}{\bm{\theta}}
\newcommand{\thetah}{\hat{\bm{\theta}}}
\newcommand{\lamv}{\bm{\lambda}}
\newcommand{\fx}{f(\xv)}
\newcommand{\fxt}{f(\xv \,|\, \bm{\theta})}
\newcommand{\fxi}{f\!\left(\xv^{(i)}\right)}
\newcommand{\fh}{\hat{f}}
\newcommand{\fxh}{\hat{f}(\xv)}
\newcommand{\yh}{\hat{y}}
\newcommand{\yih}{\hat{y}^{(i)}}
% \xi clashes with the Greek letter \xi; write \xv^{(i)} inline instead.
\newcommand{\yi}[1][i]{y^{(#1)}}
\newcommand{\xyi}[1][i]{\left(\xv^{(#1)},\, y^{(#1)}\right)}
\newcommand{\thx}{\bm{\theta}^\top \xv}
\newcommand{\sumin}{\sum_{i=1}^{n}}
\newcommand{\sumjp}{\sum_{j=1}^{p}}
\newcommand{\Lxy}{L\!\left(y, \fx\right)}
\newcommand{\Lxyi}{L\!\left(\yi, \fxi\right)}
\newcommand{\risk}{\mathcal{R}}
\newcommand{\riske}{\mathcal{R}_{\text{emp}}}
\newcommand{\risket}{\mathcal{R}_{\text{emp}}(\bm{\theta})}
\newcommand{\riskt}{\mathcal{R}(\bm{\theta})}
\newcommand{\riskr}{\mathcal{R}_{\text{reg}}}
\newcommand{\argmin}{\mathop{\mathrm{arg\,min}}}
\newcommand{\argmax}{\mathop{\mathrm{arg\,max}}}
\newcommand{\eps}{\varepsilon}

\newcolumntype{L}[1]{>{\raggedright\arraybackslash}p{#1}}

\title{Lecture 37c: Explaining TabPFN}
\subtitle{Feature effects, feature importance, and the value of training rows}
\date{}

\begin{document}

\begin{frame}\titlepage\end{frame}

\begin{frame}{One model, three explanation questions}
  \small
  Imagine a TabPFN loan-risk predictor and one applicant. We can ask:
  \vskip 8pt
  \begin{enumerate}\itemsep 7pt
    \item \textbf{Feature effect:} how does this applicant's predicted risk change
      if we vary credit utilization?
    \item \textbf{Feature importance:} how much worse is prediction on held-out
      applicants if the learner never gets that column?
    \item \textbf{Data value:} which labeled applicants help when only a subset
      fits in the context?
  \end{enumerate}
  \vskip 8pt
  Rundel, Kobialka, von Crailsheim, Feurer, Nagler and R\"ugamer (2024)
  adapt established methods to answer these questions efficiently for the
  \textbf{original, small-context TabPFN}.
  \vskip 4pt
  {\scriptsize Source: Rundel et al., \emph{Interpretable Machine Learning for TabPFN}, arXiv:2403.10923.}
\end{frame}

\begin{frame}{Outline}\tableofcontents\end{frame}

\section{What is held fixed?}
\begin{frame}[plain]
  \vfill\centering
  {\Large\color{popblue}\textbf{What is held fixed?}}\\[9pt]
  {\small An explanation changes when the context changes.}
  \vfill
\end{frame}

\begin{frame}{The prediction depends on two inputs}
  \small
  Write the in-context predictor as
  \[\hat p(y^*\mid x^*,D),\qquad D=\{(x_i,y_i)\}_{i=1}^{n}.\]
  \vskip 7pt
  \begin{columns}[T]
    \column{0.49\textwidth}
    \fcolorbox{armblue}{armblue!8}{\parbox{0.94\linewidth}{\small
      \textbf{Change the query $x^*$.}\par
      An individual conditional expectation curve varies one feature for
      the same applicant while keeping $D$ fixed.}}
    \column{0.49\textwidth}
    \fcolorbox{armred}{armred!8}{\parbox{0.94\linewidth}{\small
      \textbf{Change the context $D$.}\par
      Leave-one-out removes a training row; feature-subset methods remove a
      column from both $D$ and $x^*$.}}
  \end{columns}
  \vskip 10pt
  The network weights stay fixed in both cases. TabPFN can evaluate the new
  supervised problem with another forward pass.
  \vskip 4pt
  {\scriptsize Source: Rundel et al. (2024), Sections 2--3.}
\end{frame}

\section{Feature effects and importance}
\begin{frame}[plain]
  \vfill\centering
  {\Large\color{popblue}\textbf{Feature effects and importance}}\\[9pt]
  {\small A changed prediction and a lost predictive signal are different measurements.}
  \vfill
\end{frame}

\begin{frame}{Three ways to plot a feature's effect}
  \small
  Let $z$ be a grid of credit-utilization values.
  \vskip 8pt
  \begin{tabular}{L{0.29\linewidth}L{0.62\linewidth}}
    \textbf{Method} & \textbf{Question answered}\\[3pt]\hline
    Individual conditional expectation (ICE) & For one applicant, what is
      $\hat p(y^*\mid x^*_{j}\!\leftarrow\! z,D)$ at each $z$?\\[5pt]
    Partial dependence (PD) & Average those altered predictions over many applicants.\\[5pt]
    Accumulated local effects (ALE) & Average changes over observed local ranges,
      reducing extrapolation into empty combinations.
  \end{tabular}
  \vskip 9pt
  ICE and PD may query implausible combinations, such as very low income and a huge
  loan. ALE limits that problem but still describes \textbf{predictions}, not interventions.
  \vskip 4pt
  {\scriptsize Source: Rundel et al. (2024), Section 3.1; Apley and Zhu (2020) for ALE.}
\end{frame}

\begin{frame}{Batch the altered queries}
  \small
  Ordinary PD implementations may call the predictor once per grid point.
  With TabPFN, place all altered applicants into one inference array:
  \vskip 10pt
  \begin{center}
  \begin{tikzpicture}[>=stealth, every node/.style={font=\small}]
    \node[draw=popblue,fill=armblue!8,rounded corners,minimum width=2.7cm,minimum height=.85cm] (a) at (0,0) {fixed context $D$};
    \node[draw=armorange,fill=armorange!12,rounded corners,minimum width=3.1cm,minimum height=.85cm] (b) at (4.2,0) {all grid queries};
    \node[draw=paramgreen,fill=paramgreen!8,rounded corners,minimum width=2.8cm,minimum height=.85cm] (c) at (8.5,0) {one pass};
    \draw[->,thick] (a) -- (b); \draw[->,thick] (b) -- (c);
  \end{tikzpicture}
  \end{center}
  \vskip 10pt
  The paper gives the schematic connection-count comparison
  \[O((n_{\rm train}^{2}+n_{\rm train}n_{\rm inf})G)
    \quad\hbox{versus}\quad
    O(n_{\rm train}^{2}+n_{\rm train}n_{\rm inf}G).\]
  $G$ is the grid size. The saving comes from processing the training context once.
  \vskip 3pt
  {\scriptsize Source: Rundel et al. (2024), Section 3.1.}
\end{frame}

\begin{frame}{Does a column help on held-out rows?}
  \small
  \textbf{Leave-One-Covariate-Out (LOCO)} compares predictive risk with and
  without feature $j$. Both versions must use the \emph{same} validation rows:
  \[\mathrm{LOCO}_{j}=R_{\rm val}(D_{-j})-R_{\rm val}(D).\]
  \vskip 7pt
  \begin{center}
  \begin{tabular}{lcc}
    & all columns & without credit utilization\\[3pt]\hline
    Illustrative validation log loss & 0.22 & 0.29\\[3pt]
  \end{tabular}
  \end{center}
  \vskip 7pt
  Thus $\mathrm{LOCO}_{j}=0.29-0.22=\boxed{0.07}$. The column provides useful
  predictive information in this evaluation. A near-zero value can also mean
  another column carries the same information.
  \vskip 4pt
  {\scriptsize Method: Rundel et al. (2024), Section 3.2. Numbers are illustrative, not paper results.}
\end{frame}

\section{Shapley values through in-context learning}
\begin{frame}[plain]
  \vfill\centering
  {\Large\color{popblue}\textbf{Shapley values}}\\[9pt]
  {\small A feature's contribution depends on which other features are present.}
  \vfill
\end{frame}

\begin{frame}{A coalition is a complete smaller learning problem}
  \small
  For a subset of features $S$, define a value $v(S)$ from the model trained
  \emph{in context} on only columns in $S$ and queried on those same columns.
  \vskip 8pt
  \begin{center}
  \begin{tikzpicture}[>=stealth, every node/.style={font=\small}]
    \node[draw=armblue,fill=armblue!8,rounded corners,minimum width=3.3cm,minimum height=.9cm] (a) at (0,0) {$D$: all columns};
    \node[draw=armorange,fill=armorange!12,rounded corners,minimum width=3.3cm,minimum height=.9cm] (b) at (4.5,0) {$D_S$: columns in $S$};
    \node[draw=paramgreen,fill=paramgreen!8,rounded corners,minimum width=3.3cm,minimum height=.9cm] (c) at (9,0) {$\hat p(y^*\mid x^*_S,D_S)$};
    \draw[->,thick] (a) -- (b);
    \draw[->,thick] (b) -- (c);
  \end{tikzpicture}
  \end{center}
  \vskip 10pt
  The \textbf{feature-subset learner} is evaluated directly. We do not fill
  missing features in the query with sampled values while retaining the full
  model. This differs from ordinary Kernel SHapley Additive exPlanations
  (Kernel SHAP).
  \vskip 4pt
  {\scriptsize Source: Rundel et al. (2024), Section 3.3. Schematic.}
\end{frame}

\begin{frame}{Shapley values average the marginal contribution}
  \small
  With $p$ features, feature $j$ gets
  \[\phi_j=\sum_{S\subseteq F\setminus\{j\}}
    \frac{|S|!(p-|S|-1)!}{p!}\,[v(S\cup\{j\})-v(S)].\]
  \vskip 6pt
  For two features, the weights are both $1/2$. In an illustrative game:
  \vskip 5pt
  \begin{center}
    $v(\varnothing)=0.40,\quad v(A)=0.70,\quad
      v(B)=0.50,\quad v(A,B)=0.80.$
  \end{center}
  \vskip 5pt
  \[\phi_A=\tfrac12(0.70-0.40)+\tfrac12(0.80-0.50)=\boxed{0.30},\]
  \[\phi_B=\tfrac12(0.50-0.40)+\tfrac12(0.80-0.70)=\boxed{0.10}.\]
  The parts sum to $0.40=v(A,B)-v(\varnothing)$.
  \vskip 3pt
  {\scriptsize Definition: Shapley (1953); adaptation: Rundel et al. (2024). Values are illustrative.}
\end{frame}

\begin{frame}{"Exact retraining" does not mean exact Shapley values}
  \small
  The paper uses \textbf{exact retraining} for one particular operation:
  obtain $v(S)$ directly by passing the restricted table $D_S$ through TabPFN.
  There is no parameter optimization on $D_S$.
  \vskip 9pt
  \begin{columns}[T]
    \column{0.48\textwidth}
    \fcolorbox{paramgreen}{paramgreen!8}{\parbox{0.94\linewidth}{\small
      \textbf{Exact for a chosen coalition}\par
      Evaluate the restricted learner directly. No missing-feature imputation
      approximates that learner.}}
    \column{0.48\textwidth}
    \fcolorbox{armorange}{armorange!12}{\parbox{0.94\linewidth}{\small
      \textbf{Still approximate overall}\par
      Kernel SHAP samples $M$ feature coalitions instead of enumerating all
      $2^p$ coalitions.}}
  \end{columns}
  \vskip 10pt
  The paper reports lower attribution error for this approach at a fixed
  connection budget on its three small classification datasets, with some
  exceptions at low budgets.
  \vskip 4pt
  {\scriptsize Source: Rundel et al. (2024), Sections 3.3 and 4.2.}
\end{frame}

\section{Which training rows help?}
\begin{frame}[plain]
  \vfill\centering
  {\Large\color{popblue}\textbf{Which training rows help?}}\\[9pt]
  {\small The context is part of the predictor, so its rows can be valued too.}
  \vfill
\end{frame}

\begin{frame}{Leave-one-out measures one row's effect}
  \small
  For training row $i$, compare the held-out loss with the full context and
  with row $i$ removed:
  \[\mathrm{LOO}_i=R_{\rm val}(D_{-i})-R_{\rm val}(D).\]
  \vskip 8pt
  \begin{itemize}\itemsep 7pt
    \item Positive score: removing this row hurts prediction on the validation set.
    \item Negative score: removing it helps on that validation set.
    \item A large score may reflect a rare useful case, label noise, or validation-set choice.
  \end{itemize}
  \vskip 8pt
  One row removal is one changed context and one forward pass. The paper also
  explores gradients with respect to training observations, a different notion
  of sensitivity.
  \vskip 4pt
  {\scriptsize Source: Rundel et al. (2024), Sections 3.4--3.5.}
\end{frame}

\begin{frame}{Data Shapley selects a small context from a large pool}
  \small
  The 2024 model accepted roughly 1,024 context rows. Suppose the available
  labeled pool is larger. A random 512-row subset wastes some informative rows.
  \vskip 7pt
  \begin{enumerate}\itemsep 5pt
    \item Sample many candidate subsets of training rows.
    \item Evaluate each subset on a separate validation set.
    \item Attribute the change in validation risk to individual rows.
    \item Choose 512 rows with favorable estimated contributions.
  \end{enumerate}
  \vskip 8pt
  In the paper, 3,072 available rows are reduced to 512. Across three selected
  OpenML datasets and five runs, the optimized context beats random selection
  in reported test ROC AUC. The procedure takes \textbf{nearly 10,000 forward passes}.
  \vskip 4pt
  {\scriptsize Source: Rundel et al. (2024), Sections 3.4 and 4.3.}
\end{frame}

\begin{frame}{The package also covers other explanation tasks}
  \small
  Rundel et al. provide more than the four methods we worked through:
  \vskip 8pt
  \begin{tabular}{L{0.32\linewidth}L{0.58\linewidth}}
    \textbf{Method in the paper} & \textbf{Question or role}\\[3pt]\hline
    Shapley Additive Global Importances (SAGE) & Attribute predictive risk to
      features across a dataset; uses the subset-evaluation idea.\\[5pt]
    Sensitivity analysis & Use gradients of predictions or risk to score local
      feature or training-row sensitivity.\\[5pt]
    Conformal prediction & Wrap predictions to form uncertainty sets.\\[5pt]
    Counterfactual explanations and decision curves & Explore alternative inputs
      and assess decisions; package integrations rather than the main new result.
  \end{tabular}
  \vskip 9pt
  An explanation method is useful only after its \textbf{target quantity} is clear.
  Sensitivity, predictive risk, and decision value answer different questions.
  \vskip 4pt
  {\scriptsize Source: Rundel et al. (2024), Table 1 and Sections 3.5--3.6.}
\end{frame}

\section{Interpretation limits}
\begin{frame}[plain]
  \vfill\centering
  {\Large\color{popblue}\textbf{Interpretation limits}}\\[9pt]
  {\small A useful explanation still needs the right question and evaluation data.}
  \vfill
\end{frame}

\begin{frame}{Predictive explanations are not causal effects}
  \small
  In the loan example, a PD curve may rise as utilization rises. That means
  \emph{the predictor's output} changes when we edit that input. It does not
  establish that changing a person's utilization would change their actual outcome.
  \vskip 10pt
  \begin{columns}[T]
    \column{0.49\textwidth}
    \fcolorbox{armblue}{armblue!8}{\parbox{0.94\linewidth}{\small
      \textbf{Prediction question}\par
      What does $\hat p(y\mid x,D)$ do under an input or context change?}}
    \column{0.49\textwidth}
    \fcolorbox{armred}{armred!8}{\parbox{0.94\linewidth}{\small
      \textbf{Causal question}\par
      What would happen to the outcome under a real intervention? That needs
      causal assumptions and a different design.}}
  \end{columns}
  \vskip 10pt
  Correlated columns make attribution ambiguous: LOCO can be small because
  another feature substitutes for the removed one.
  \vskip 4pt
  {\scriptsize Methods: Rundel et al. (2024). Causal distinction follows the estimands defined here.}
\end{frame}

\begin{frame}{The paper studied the 2024 model}
  \small
  The authors built \texttt{tabpfn\_iml} for the original TabPFN, whose practical
  context was about 1,024 rows. The 2.5, 3 and 3.5 releases came later.
  \vskip 9pt
  \begin{tabular}{L{0.31\linewidth}L{0.59\linewidth}}
    \textbf{Portable idea} & A context change can evaluate a new restricted learner
      without optimizing model weights.\\[7pt]
    \textbf{Needs rechecking} & Runtime, memory, input encodings, attribution
      stability, and performance on later checkpoints.\\[7pt]
    \textbf{Evaluation rule} & Keep a held-out set for LOCO and row valuation;
      inspect dependence between features before giving a causal reading.
  \end{tabular}
  \vskip 8pt
  The paper and implementation are a method demonstration, not a validation
  of every explanation method on TabPFN-3.5.
  \vskip 4pt
  {\scriptsize Source: Rundel et al. (2024), arXiv:2403.10923; later versions: Prior Labs (2025, 2026).}
\end{frame}

\section*{Wrap-up}
\begin{frame}{Recap}
  \small
  \begin{enumerate}\itemsep 7pt
    \item ICE, PD and ALE vary \textbf{query features}; batching saves repeated
      context processing.
    \item LOCO removes a \textbf{column} and measures the held-out loss difference.
    \item The paper evaluates a Shapley coalition by restricting \textbf{both}
      context and query; sampled coalitions still approximate the full sum.
    \item LOO and Data Shapley value \textbf{training rows}. Better context
      selection can cost many forward passes.
  \end{enumerate}
  \vskip 8pt
  \fcolorbox{paramgreen}{paramgreen!8}{\parbox{0.96\linewidth}{\small
    \textbf{Next:} choose the question first. A local feature effect, global
    predictive importance, and training-row value are different explanations.}}
  \vskip 7pt
  {\scriptsize Paper: arxiv.org/abs/2403.10923 \quad Code: github.com/david-rundel/tabpfn\_iml}
\end{frame}

\end{document}

% Provenance:
%   Rundel, Kobialka, von Crailsheim, Feurer, Nagler and Rugamer (2024),
%   Interpretable Machine Learning for TabPFN, arXiv:2403.10923.
%   Local PDF: papers/2403.10923_interpretable-machine-learning-for-tabpfn.pdf.
%   All worked numbers are labeled illustrative. TikZ figures are schematics.
%   The paper evaluates the original TabPFN, not later 2.5/3/3.5 checkpoints.
%   Forward pointer: apply the methods to a chosen model and held-out split.
